\section{The Quotient Graph}

Let
\[
Q_{15} = G_{60}/V_4
\]
denote the quotient graph obtained from the action described in the
previous section.

The vertices of $Q_{15}$ correspond to the $15$ orbits of the
$V_4$-action on $V(G_{60})$.

\begin{proposition}
The quotient graph $Q_{15}$ has
\[
|V(Q_{15})| = 15, \qquad |E(Q_{15})| = 30 .
\]
\end{proposition}

\begin{proof}
Each orbit contains four vertices of $G_{60}$, so the quotient graph
has $60/4=15$ vertices.

Since the action preserves adjacency and $G_{60}$ is $4$-regular, the
quotient graph remains $4$-regular.  Thus the number of edges is
\[
|E(Q_{15})| = \frac{1}{2}\cdot 15\cdot 4 = 30 .
\]
\end{proof}

\begin{proposition}
The shell profile of $Q_{15}$ is
\[
(1,4,8,2).
\]
\end{proposition}

In the next section we identify this graph with a natural graph
arising from the geometry of the dodecahedron.
